\ifdefined\ishandout
  \documentclass[handout]{beamer}
\else
  \documentclass{beamer}
\fi

\usetheme{boxes}
\definecolor{beamer@structure@color}{rgb}{0,0,0}

\usecolortheme{structure}

\setbeamertemplate{footline}[frame number]
\setbeamertemplate{frametitle}{\color{black}
\def\myhrulefill{\leavevmode\leaders\hrule height 1pt\hfill\kern 0pt}
\headingfont\insertframetitle\par\vskip-8pt\myhrulefill}

\newcommand{\personality}[1]{{\bf #1}}
\usepackage{cancel}

\newcommand{\NN}{\mathbb{N}}
\newcommand{\ZZ}{\mathbb{Z}}
\newcommand{\QQ}{\mathbb{Q}}
\newcommand{\RR}{\mathbb{R}}
\newcommand{\CC}{\mathbb{C}}
\newcommand{\FF}{\mathbb{F}}
\renewcommand{\AA}{\mathbb{A}}
\newcommand{\isom}{\cong}
\renewcommand{\O}{\mathcal{O}}
\newcommand{\dfn}{\mathrel{\mathop:}=}
\DeclareMathOperator{\ord}{ord}
\DeclareMathOperator{\rk}{rk}
\DeclareMathOperator{\Spec}{Spec}
\renewcommand{\Re}{\mathop{\mathrm{Re}}}
\DeclareMathOperator{\Hom}{Hom}
\DeclareMathOperator{\Gal}{Gal}
\newcommand{\RHom}{R\!\Hom}
\newcommand{\categ}[1]{\text{\bf #1}}
\newcommand{\quiso}{\simeq}

\setbeamertemplate{navigation symbols}{}

\usepackage{array}
\newcolumntype{x}[1]{>{\centering\hspace{0pt}}p{#1}}
\definecolor{shadecolor}{rgb}{0.89,0.89,0.89}
\usepackage{colortbl}

\newcommand{\term}{\textbf}

\usepackage{mathspec}
\setsansfont[BoldFont={PT Sans Bold}, ItalicFont={PT Sans Italic}]{PT Sans}
\setmathrm[BoldFont={PT Sans Bold}, ItalicFont={PT Sans Italic}]{PT Sans}
\newfontfamily\headingfont[]{PT Sans Bold}

\usepackage[all]{xy}

\usepackage{tikz-cd}
\usetikzlibrary{arrows}
\usetikzlibrary{calc}

\begin{document}

% % % % % % % % % % % % % % % % % % % % % % % % % % % % % % % % % % % % % % % % % % % % %

\begin{frame}[noframenumbering]
  \headingfont
  \begin{center}
    {\huge Zeta-values of arithmetic schemes at negative integers and Weil-étale
      cohomology

    }

    \vspace{2em}

    {\large Alexey Beshenov}

    \vspace{0.20em}

    Universit\'e de Bordeaux / Universiteit Leiden

    \vspace{1em}

    {\large Advised by Baptiste Morin and Bas Edixhoven}

    \vspace{1em}

    {\tiny 4 December 2017, Leiden}

    \vspace{5em}

    \raisebox{+0.125cm}{\includegraphics[width=2.5cm]{../u-bordeaux.pdf}}\hspace{0.9cm}
    \raisebox{+0.12cm}{\includegraphics[width=2.25cm]{../leiden.pdf}}\hspace{0.9cm}
    \includegraphics[width=2.5cm]{../algant.mps}
  \end{center}
\end{frame}

% % % % % % % % % % % % % % % % % % % % % % % % % % % % % % % % % % % % % % % % % % % % %

\begin{frame}
  \frametitle{Outline}

  \begin{itemize}
  \item<2-> Motivation

  \item<3-> Motivic cohomology

  \item<4-> Weil-étale complexes

  \item<5-> Regulator

  \item<6-> Main conjecture

  \item<7-> Stability properties

  \item[*]<8-> Definition of Weil-étale complexes
  \end{itemize}

\end{frame}

% % % % % % % % % % % % % % % % % % % % % % % % % % % % % % % % % % % % % % % % % % % % %

\begin{frame}[plain]
  \headingfont

  \vspace{\fill}

  \begin{center}
    {\huge Motivation}
  \end{center}

  \vspace{\fill}
\end{frame}

% % % % % % % % % % % % % % % % % % % % % % % % % % % % % % % % % % % % % % % % % % % % %

\begin{frame}
  \frametitle{Motivation: zeta values}

  \begin{itemize}
  \item<2-> \term{Zeta function} of an \term{arithmetic scheme}:

    \[ \begin{tikzcd}
        X \ar{d}[swap]{\begin{array}{r}\text{separated}\\\text{of finite type}\end{array}} \\
        \Spec \ZZ
      \end{tikzcd}
      \!\!\!\!\!\!\! \leadsto ~
      \zeta (X, s) \dfn \prod_{x\in X_0} \frac{1}{1 - N (x)^{-s}}
      \quad (\Re s > \dim X) \]

  \item<3-> \textbf{Conjecture}: meromorphic continuation to $\CC$.

  \item<4-> Two quantities for each $n\in \ZZ$:

    \onslide<5->{$d_n \dfn \ord_{s = n} \zeta (X,s) \dfn$
      \term{vanishing order} at $s=n$,}

    \onslide<6->{\term{special value}
      $\zeta^* (X,n) \dfn \lim\limits_{s\to n} (s-n)^{-d_n} \, \zeta (X,s)$.}
  \end{itemize}
\end{frame}

% % % % % % % % % % % % % % % % % % % % % % % % % % % % % % % % % % % % % % % % % % % % %

\begin{frame}
  \frametitle{Motivation:\\
    cohomological interpretation of zeta values}

  \begin{itemize}
  \item<2-> May be traced to \personality{Dirichlet} (1839):
    $$\zeta_F^* (0) = \zeta^* (\Spec \O_F, 0) = -\frac{h_F}{\#\mu_F}\,R_F.$$

  \item<3-> Actually,
    $$\cdots = -\frac{\# H^0 (\Spec \O_F, \ZZ^c (0))}{\# H^{-1} (\Spec \O_F, \ZZ^c (0))_\text{\it tors}}\,R_F,$$
    where $\ZZ^c (0) \isom \mathbb{G}_m [1]$ (to be explained in a moment).

  \item<4->[*] First results about finite generation of motivic cohomology are
    from the XIX century :--)

    $H^0 (\Spec \O_F, \ZZ^c (0))$ is finite;\\
    $H^{-1} (\Spec \O_F, \ZZ^c (0))$ is of rank $r_1 + r_2 - 1.$
  \end{itemize}
\end{frame}

% % % % % % % % % % % % % % % % % % % % % % % % % % % % % % % % % % % % % % % % % % % % %

\begin{frame}
  \frametitle{Motivation:\\
    Weil-étale cohomology (very brief history)}

  \begin{itemize}
  \item<2-> \personality{Lichtenbaum}: $d_n$ and $\zeta^* (X,n)$ conjecturally
    come from ``\term{Weil-étale cohomology}''.

  \item<3-> \personality{Lichtenbaum}, 2005: first studies for $X_{/\FF_q}$.

  \item<4-> \personality{Geisser}, 2004: further results for $X_{/\FF_q}$.

  \item<5-> \personality{Lichtenbaum}, 2009: $X = \Spec \O_F$, $n = 0$.

  \item<6-> \personality{Morin}, 2014: $X_{/\ZZ}$ proper, regular, $n = 0$.

  \item<7-> \personality{Flach} and \personality{Morin}, 2016: $X_{/\ZZ}$
    proper, regular, $n \in \ZZ$.
  \end{itemize}
\end{frame}

% % % % % % % % % % % % % % % % % % % % % % % % % % % % % % % % % % % % % % % % % % % % %

\begin{frame}
  \frametitle{Motivation: my thesis}

  \begin{itemize}
  \item<2-> \textbf{Goal}: construct Weil-étale cohomology following\\
    Flach and Morin...

  \item<3-> For \emph{any} arithmetic scheme $X$ (harder).

  \item<4-> For $n < 0$ (easier).
  \end{itemize}

  \begin{center}\onslide<5->
    {\large From now on, $n$ is strictly negative!}
  \end{center}

\end{frame}

% % % % % % % % % % % % % % % % % % % % % % % % % % % % % % % % % % % % % % % % % % % % %

\begin{frame}[plain]
  \headingfont

  \vspace{\fill}

  \begin{center}
    {\huge Motivic cohomology}
  \end{center}

  \vspace{\fill}
\end{frame}

% % % % % % % % % % % % % % % % % % % % % % % % % % % % % % % % % % % % % % % % % % % % %

\begin{frame}
  \frametitle{Motivic cohomology}

  \begin{itemize}
  \item<2-> \term{Bloch}, 1986: \term{cycle complexes} $\ZZ (n)$,\\
    \term{higher Chow groups} $CH^i (X,p) = H^{2i - p} (X, \ZZ (i))$.

  \item<3-> \term{Geisser}, 2010: \term{dualizing complexes}.
    $$(X,n) \rightsquigarrow \ZZ^c (n) \text{, complex of sheaves on }X_\text{\it ét}.$$

  \item<4-> $\ZZ^c (n) = \ZZ (d-n) [2d], \quad d = \dim X$.

  \item<5->[*] For $X_{/k}$ smooth, RHS $\isom$ Voevodsky's motivic complex\\
    (not our case).
  \end{itemize}
\end{frame}

% % % % % % % % % % % % % % % % % % % % % % % % % % % % % % % % % % % % % % % % % % % % %

\begin{frame}
  \frametitle{Motivic cohomology:\\
    Borel--Moore behavior of $\ZZ^c (n)$}

  \begin{itemize}
  \item<2-> \term{Borel--Moore homology} for locally compact spaces.

  \item<3-> \term{Localization triangles} for $U \hookrightarrow X \leftarrow Z$:
    $$R\Gamma_{BM} (Z, \ZZ) \to
    R\Gamma_{BM} (X, \ZZ) \to
    R\Gamma_{BM} (U, \ZZ) \to \cdots [1]$$

  \item<4-> \term{Verdier}:
    $R\Gamma_{BM} (X,\ZZ) = R\Gamma (X, p^! \underline{\ZZ}) \isom
    \RHom (R\Gamma_c (X,\ZZ), \ZZ)$, where $p\colon X\to \ast$.

  \item<5-> Dual to cohomology with compact support:
    $$R\Gamma_c (U,\ZZ) \to R\Gamma_c (X,\ZZ) \to R\Gamma_c (Z,\ZZ) \to \cdots [1]$$

  \item<6-> \term{Geisser}, 2010: for schemes,
    $$R\Gamma (Z_\text{\it ét}, \ZZ^c (n)) \to
    R\Gamma (X_\text{\it ét}, \ZZ^c (n)) \to
    R\Gamma (U_\text{\it ét}, \ZZ^c (n)) \to \cdots [1]$$
  \end{itemize}
\end{frame}

% % % % % % % % % % % % % % % % % % % % % % % % % % % % % % % % % % % % % % % % % % % % %

\begin{frame}
  \frametitle{Motivic cohomology:\\
    assumption on finite generation}

  \begin{itemize}
  \item<2-> \term{Conjecture} $\mathbf{L}^c (X_\text{\it ét}, n)$:

    \emph{$H^i (X_\text{\it ét}, \ZZ^c (n))$
      are finitely generated for all $i\in\ZZ$}.

  \item<3->[*] As always here, $n < 0$.
  \end{itemize}
\end{frame}

% % % % % % % % % % % % % % % % % % % % % % % % % % % % % % % % % % % % % % % % % % % % %

\begin{frame}[plain]
  \headingfont

  \vspace{\fill}

  \begin{center}
    {\huge Weil-étale complexes}
  \end{center}

  \vspace{\fill}
\end{frame}

% % % % % % % % % % % % % % % % % % % % % % % % % % % % % % % % % % % % % % % % % % % % %

\begin{frame}
  \frametitle{Weil-étale complexes}

  \onslide<2->{\textbf{Theorem}.}

  \begin{itemize}
  \item<3-> For an arithmetic scheme $X$ and $n < 0$, assume
    $\mathbf{L}^c (X_\text{\it ét}, n)$.

  \item<4-> We may construct a perfect complex of abelian groups
    $$R\Gamma_\text{\it W,c} (X, \ZZ (n))$$
    (well-defined up to an iso in $\categ{D} (\ZZ)$).

  \item<5-> \term{perfect} = $H^i_\text{\it W,c} (X, \ZZ (n))$ are finitely generated;\\
    zero for almost all $i$.
  \end{itemize}
\end{frame}

% % % % % % % % % % % % % % % % % % % % % % % % % % % % % % % % % % % % % % % % % % % % %

\begin{frame}
  \frametitle{Weil-étale complexes: splitting}

  \begin{itemize}
  \item<2-> $R\Gamma_\text{\it W,c} (X,\ZZ(n)) \otimes \RR \isom
    \RHom (R\Gamma (X_\text{\it ét}, \ZZ^c (n)), \RR) [-1] \oplus
    R\Gamma_c (G_\RR, X (\CC), (2\pi i)^n\,\RR) [-1]$.

  \item<3-> $G_\RR \dfn \Gal (\CC/\RR)$ acts on
    $R\Gamma_c (X (\CC), (2\pi i)^n\,A)$ $\Rightarrow$ $G_\RR$-equivariant
    cohomology with compact support:
    $$R\Gamma_c (G_\RR, X (\CC), (2\pi i)^n\,A) \dfn
    R\Gamma (G_\RR, R\Gamma_c (X (\CC), (2\pi i)^n\,A)).$$

  \item<4-> Grothendieck spectral sequence:
    $$E_2^{pq} = H^p (G_\RR, H^q_c (X (\CC), (2\pi i)^n\,A)) \Longrightarrow
    H^{p+q}_c (G_\RR, X (\CC), (2\pi i)^n\,A).$$

  \item<5-> Note:
    $$H_c^i (G_\RR, X (\CC), (2\pi i)^n\,\RR) \isom H_c^i (X (\CC), (2\pi i)^n\,\RR)^{G_\RR}$$
    ($H^{>0} (G_\RR, \ldots)$ is always $2$-torsion!)
  \end{itemize}
\end{frame}

% % % % % % % % % % % % % % % % % % % % % % % % % % % % % % % % % % % % % % % % % % % % %

\begin{frame}[plain]
  \headingfont

  \vspace{\fill}

  \begin{center}
    {\huge Regulator}
  \end{center}

  \vspace{\fill}
\end{frame}

% % % % % % % % % % % % % % % % % % % % % % % % % % % % % % % % % % % % % % % % % % % % %

\begin{frame}
  \frametitle{Regulator}

  \begin{itemize}
  \item<2-> A nice construction due to\\
    \personality{Kerr, Lewis, and Müller-Stach} (2006).

  \item<3-> Usually:
    $$\begin{array}{c}
        \text{motivic cohomology}\\
        \text{(higher Chow groups)}
      \end{array}\text{of }X \xrightarrow{Reg} \text{Deligne (co)homology of }X(\CC).$$

  \item<4-> In our case $n < 0$: a morphism in $\categ{D} (\ZZ)$
    $$Reg\colon R\Gamma(X_\text{\it ét},\ZZ^c (n)) \to
    R\Gamma_\text{\it BM} (G_\RR,X(\CC),(2\pi i)^n\,\RR)[1].$$

  \item<5-> \textbf{Drawback}: need to assume that $X_\CC$ is smooth,
    quasi-projective.

  \item<6-> \textbf{Conjecture} $\mathbf{B}(X, n)$: the $\RR$-dual is a
    quasi-isomorphism
    $$Reg^\vee\colon R\Gamma_c (G_\RR, X(\CC), (2\pi i)^n\,\RR)[-1] \xrightarrow{\quiso}
    \RHom(R\Gamma(X_\text{\it ét}, \ZZ^c (n)), \RR)$$
  \end{itemize}
\end{frame}

% % % % % % % % % % % % % % % % % % % % % % % % % % % % % % % % % % % % % % % % % % % % %

\begin{frame}
  \frametitle{Morphism $\smile\theta$}

  \begin{itemize}
  \item<2-> Define
    $$R\Gamma_\text{\it W,c} (X, \ZZ(n)) \otimes \RR \xrightarrow{\smile\theta}
    R\Gamma_\text{\it W,c} (X, \ZZ(n)) \otimes \RR [1]$$

    \[ \begin{tikzcd}[row sep=small]
        R\Gamma_\text{\it W,c} (X, \ZZ (n)) \otimes \RR \ar{d}{\isom}[swap]{\text{splitting}} \\
        \RHom (R\Gamma (X_\text{\it ét}, \ZZ^c (n)), \RR) [-1] ~\oplus~ R\Gamma_c (G_\RR, X (\CC), (2\pi i)^n\,\RR) [-1] \ar[twoheadrightarrow]{d} \\
        R\Gamma_c (G_\RR, X (\CC), (2\pi i)^n\,\RR) [-1] \ar{d}{Reg^\vee}[swap]{\text{conjecturally }\isom}\\
        \RHom (R\Gamma (X_\text{\it ét}, \ZZ^c (n)), \RR) \ar[rightarrowtail]{d} \\
        \RHom (R\Gamma (X_\text{\it ét}, \ZZ^c (n)), \RR) ~\oplus~ R\Gamma_c (G_\RR, X (\CC), (2\pi i)^n\,\RR) \ar{d}{\isom}[swap]{\text{splitting}} \\
        R\Gamma_\text{\it W,c} (X, \ZZ (n)) \otimes \RR [1]
      \end{tikzcd} \]
  \end{itemize}
\end{frame}

% % % % % % % % % % % % % % % % % % % % % % % % % % % % % % % % % % % % % % % % % % % % %

\begin{frame}
  \frametitle{Trivialization map $\lambda$}

  \begin{itemize}
  \item<2-> Assume $\mathbf{L}^c (X_\text{\it ét},n)$ and
    $\mathbf{B} (X,n)$. Then

    \[ \tag{*} \cdots \to H_\text{\it W,c}^i (X, \ZZ(n)) \otimes \RR
      \xrightarrow{\smile\theta}
      H_\text{\it W,c}^{i+1} (X, \ZZ(n)) \otimes \RR \to \cdots \]

    is an acyclic complex of finite dimensional vector spaces.

  \item<3-> \personality{Knudsen, Mumford}, 1976: determinants of complexes.
    $$C^\bullet \text{, perfect object in }\categ{D} (\ZZ) \leadsto
    \det\nolimits_\ZZ C^\bullet\text{, }\ZZ\text{-module of rank }1.$$

  \item<4-> \textbf{Properties of determinants} $\Rightarrow$ (*)
    induces a \emph{canonical} iso
    $$\lambda\colon \RR \xrightarrow{\isom}
    (\det\nolimits_\ZZ R\Gamma_\text{\it W,c} (X,\ZZ(n)) \otimes \RR.$$

  \item<5-> $R\Gamma_\text{\it W,c} (X,\ZZ(n))$ is constructed up to iso in
    $\categ{D} (\ZZ)$, but $\det\nolimits_\ZZ R\Gamma_\text{\it W,c} (X,\ZZ(n))$
    is \emph{canonically} defined and \emph{canonically} seen as a lattice in
    $\RR$.
  \end{itemize}
\end{frame}

% % % % % % % % % % % % % % % % % % % % % % % % % % % % % % % % % % % % % % % % % % % % %

\begin{frame}[plain]
  \headingfont

  \vspace{\fill}

  \begin{center}
    {\huge Main conjecture}
  \end{center}

  \vspace{\fill}
\end{frame}

% % % % % % % % % % % % % % % % % % % % % % % % % % % % % % % % % % % % % % % % % % % % %

\begin{frame}
  \frametitle{Conjecture $\mathbf{C} (X,n)$}

  \onslide<2->{For an arithmetic scheme $X$ and $n < 0$}

  \begin{enumerate}
  \item[a)]<3-> assume $\mathbf{L}^c (X_\text{\it ét}, n)$;

  \item[b)]<4-> assume that $X_\CC$ is smooth, quasi-projective;\\
    assume $\mathbf{B}(X, n)$;

  \item[c)]<5-> assume meromorphic continuation near $s = n$ for $\zeta (X, s)$.
  \end{enumerate}

  \onslide<6->{Then}

  \begin{enumerate}
  \item[1)]<7-> the vanishing order is given by
    \[ \ord_{s=n} \zeta (X,s) =
      \sum_{i\in\ZZ} (-1)^i \cdot i\cdot \rk_\ZZ H^i_\text{\it W,c} (X, \ZZ(n)). \]

  \item[2)]<8-> the special value is given up to sign by
    \[ \lambda (\zeta^* (X, n)^{-1}) \cdot \ZZ =
      \det\nolimits_\ZZ R\Gamma_\text{\it W,c} (X, \ZZ(n)). \]
  \end{enumerate}
\end{frame}

% % % % % % % % % % % % % % % % % % % % % % % % % % % % % % % % % % % % % % % % % % % % %

\begin{frame}
  \frametitle{Main conjecture:\\
    relation to other statements}

  \begin{itemize}
  \item<2-> For $X$ regular, proper,
    $\mathbf{C} (X,n) \iff \text{conjectures of \personality{Flach} and \personality{Morin} (2016)}$.

  \item<3-> \personality{Flach} and \personality{Morin} (2016):\\
    their special value conjecture $\iff$\\
    \term{Tamagawa number conjecture}\\
    (\personality{Bloch, Kato, Fontaine, Perrin-Riou}).
  \end{itemize}
\end{frame}

% % % % % % % % % % % % % % % % % % % % % % % % % % % % % % % % % % % % % % % % % % % % %

\begin{frame}
  \frametitle{Main conjecture:\\
    meaning of the vanishing order}

  \begin{itemize}
  \item<2-> What is that sum
    $\sum_{i\in\ZZ} (-1)^i \cdot i\cdot \rk_\ZZ H^i_\text{\it W,c} (X, \ZZ(n))$?

  \item<3-> Under the assumptions a), b),

    \begin{multline*}
      \ord_{s=n} \zeta (X,s) \stackrel{?}{=}
      \sum_{i\in\ZZ} (-1)^i \cdot i\cdot \rk_\ZZ H^i_\text{\it W,c} (X, \ZZ(n)) \\
      = \sum_{i\in\ZZ} (-1)^i \, \dim_\RR H^i_c (X (\CC), (2\pi i)^n\,\RR)^{G_\RR} \\
      = \chi (R\Gamma_c (G_\RR, X(\CC), (2\pi i)^n\,\RR)).
    \end{multline*}
  \end{itemize}
\end{frame}

% % % % % % % % % % % % % % % % % % % % % % % % % % % % % % % % % % % % % % % % % % % % %

\begin{frame}
  \frametitle{Toy example: number rings}

  \begin{itemize}
  \item<2-> Let $X = \Spec \O_F$, $[F : \QQ] = r_1 + 2 r_2$. The $G_\RR$-space
    $X (\CC)$:

    \[ \xymatrix@R=0.75em@C=0.5em{
        & & & & & \bullet\ar@/^0.5em/[dd] & \bullet\ar@/^0.5em/[dd] & \cdots & \bullet\ar@/^0.5em/[dd] \\
        \bullet\ar@(ul,ur) & \bullet\ar@(ul,ur) & \cdots & \bullet\ar@(ul,ur) \\
        & & & & & \bullet\ar@/^0.5em/[uu] & \bullet\ar@/^0.5em/[uu] & \cdots & \bullet\ar@/^0.5em/[uu] \\
        \ar@{}[rrr]|{r_1\text{ points}} & & & & & \ar@{}[rrr]|{2\,r_2\text{ points}} & & & \\
      } \]

  \item<3-> $R\Gamma_c (X(\CC), (2\pi i)^n\,\RR)$: a single $\RR$-vector space
    $$V \dfn ((2\pi i)^n\,\RR)^{\oplus r_1} \oplus ((2\pi i)^n\,\RR \oplus (2\pi i)^n\,\RR)^{\oplus r_2}.$$

  \item<4-> $G_\RR$-actions: $z \mapsto \overline{z}$ and
    $(z_1,z_2) \mapsto (\overline{z_2}, \overline{z_1})$.

  \item<5->
    $$\chi (R\Gamma_c (G_\RR, X (\CC), (2\pi i)^n\,\RR)) =
    \dim_\RR V^{G_\RR} = \begin{cases}
      r_2, & n \text{ odd},\\
      r_1 + r_2, & n \text{ even}.
    \end{cases}$$
  \end{itemize}
\end{frame}

% % % % % % % % % % % % % % % % % % % % % % % % % % % % % % % % % % % % % % % % % % % % %

\begin{frame}
  \frametitle{Dull example: varieties over finite fields}

  \begin{itemize}
  \item<2-> For $X_{/\FF_q}$ we have $X (\CC) = \emptyset$.

  \item<3-> \term{Weil's zeta function}:
    $$\zeta (X,s) = \exp \left(\sum_{k\ge 1} \frac{\# X (\FF_{q^k})}{k}\,q^{-ks}\right).$$

  \item<4-> For $n < 0$
    $$\ord_{s=n} \zeta (X,s) = \chi (R\Gamma_c (G_\RR, X (\CC), (2\pi i)^n\,\RR)) = 0.$$
  \end{itemize}
\end{frame}

% % % % % % % % % % % % % % % % % % % % % % % % % % % % % % % % % % % % % % % % % % % % %

\begin{frame}[plain]
  \headingfont

  \vspace{\fill}

  \begin{center}
    {\huge Okay, so how can we obtain new results?}
  \end{center}

  \vspace{\fill}
\end{frame}

% % % % % % % % % % % % % % % % % % % % % % % % % % % % % % % % % % % % % % % % % % % % %

\begin{frame}
  \frametitle{Stability properties}

  \begin{itemize}
  \item<2-> For the zeta function:

    \begin{align*}
      U  \hookrightarrow X \leftarrow Z & \leadsto \zeta (X, s) = \zeta (U, s) \cdot \zeta (Z, s),\\
      \AA^r_X \dfn \AA_\ZZ^r \times X & \leadsto \zeta (\AA^r_X, s) = \zeta (X, s - r).
    \end{align*}

  \item<3-> \textbf{Theorem}. For the conjecture:

    \begin{align*}
      U  \hookrightarrow X \leftarrow Z & \leadsto \text{two out of three } \mathbf{C} (U, n), \, \mathbf{C} (X, n), \, \mathbf{C} (Z, n) \\
                                        & \quad\quad\quad\Longrightarrow \text{the other one},\\
      \AA^r_X \dfn \AA_\ZZ^r \times X & \leadsto \mathbf{C} (\AA^r_X, n) \iff \mathbf{C} (X, n - r).
    \end{align*}

  \end{itemize}
\end{frame}

% % % % % % % % % % % % % % % % % % % % % % % % % % % % % % % % % % % % % % % % % % % % %

\begin{frame}
  \frametitle{Stability properties: proof idea}

  \begin{itemize}
  \item<2-> \textbf{Vanishing orders and localization}:\\
    additivity of $\chi$ for the triangle

    \vspace{-2em}

    \begin{multline*}
      R\Gamma_c (G_\RR, U(\CC), (2\pi i)^n\,\RR) \to
      R\Gamma_c (G_\RR, X(\CC), (2\pi i)^n\,\RR) \\
      \to R\Gamma_c (G_\RR, Z(\CC), (2\pi i)^n\,\RR) \to \cdots [1]
    \end{multline*}

  \item<3-> \textbf{Vanishing orders and affine space}:\\
    $G_\RR$-equivariant quasi-isomorphism
    $$R\Gamma_c (\CC^r \times X(\CC), (2\pi i)^n\,\RR) \quiso
    R\Gamma_c (X(\CC), (2\pi i)^{n-r}\,\RR) [-2r].$$

  \item<4-> \textbf{Special value and localization}:\\
    look also at the dual to the Borel--Moore triangle for $\ZZ^c (n)$:

    \vspace{-2em}

    \begin{multline*}
      \RHom(R\Gamma(U_\text{\it ét}, \ZZ^c (n)), \RR) \to
      \RHom(R\Gamma(X_\text{\it ét}, \ZZ^c (n)), \RR) \to\\
      \RHom(R\Gamma(Z_\text{\it ét}, \ZZ^c (n)), \RR) \to \cdots [1]
    \end{multline*}

  \item<5-> \textbf{Special value and affine space}:

    \vspace{-2em}

    \begin{multline*}
      p\colon \AA^r_X \to X \Longrightarrow R p_* \ZZ^c (n) \quiso \ZZ^c (n - r) [2r] \\
      \Longrightarrow R\Gamma (\AA^r_{X,\text{\it ét}}, \ZZ^c (n)) \quiso
      R\Gamma (X_\text{\it ét}, \ZZ^c (n - r)) [2r].
    \end{multline*}

  \end{itemize}
\end{frame}

% % % % % % % % % % % % % % % % % % % % % % % % % % % % % % % % % % % % % % % % % % % % %

\begin{frame}[plain]
  \headingfont

  \vspace{\fill}

\end{frame}

% % % % % % % % % % % % % % % % % % % % % % % % % % % % % % % % % % % % % % % % % % % % %

\begin{frame}[plain]
  \headingfont

  \vspace{\fill}

  \begin{center}
    {\huge Thank you!}
  \end{center}

  \vspace{\fill}
\end{frame}

\end{document}
